% !TEX root = ../jlt.tex
\section{The compiler's own functional language}
\label{app:fjlt:ssa}

Instantiation is partial evaluation of the dictionary-passing term at a static type, $\llbracket \mathit{concat} \rrbracket\,[\mathtt{size\_t}]\,d$ with its type and record arguments applied and $\beta$-reduced before the program runs, so the two agree on every closed term: Theorem~\ref{thm:simulation} at the binding time where the records are constants, nothing of a record surviving to run time, and parametricity, a property of the term before that evaluation, holds of every instantiation as an application of it.
For Go, Ellis et al.~\cite{ellis2022generic} prove the dictionary-passing translation correct beside monomorphisation, which Griesemer et al.~\cite{griesemer2020featherweight} proved; \Cpp{} only ever does the second, and the first carries the theorems.
What that evaluation leaves is a closed first-order term over the base types: the type abstractions are gone with their arguments, the record parameters are gone with the records the compiler resolved, every trait call is a call to one specialisation, and the only recursion left is $\mathit{iter}$ and $\mathit{fold}_F$.
The compiler then lowers that term to its intermediate representation, and the representation is in static single assignment form~\cite{cytron1991ssa}: every variable is assigned exactly once, and where two paths of control join, a $\phi$-node names the value that arrives along each.

SSA is a functional language.
Kelsey showed it to correspond to continuation-passing style~\cite{kelsey1995correspondence}, and Appel put the correspondence as a slogan~\cite{appel1998ssa}: a basic block is a local function, its $\phi$-nodes are that function's parameters, a branch is a call that passes the arguments, and dominance is lexical scope.
So a Jlt function is translated twice into a functional language, by Appendix~\ref{app:fjlt:translation} into System~F with records and iterators, and by the compiler into first-order functions with tail calls, with instantiation between the two; Rules~1 and~3 are what makes the source already look like the second.

Listing~\ref{lst:ssa} is a search under a budget as the compiler leaves it before it unrolls the loop: \cppinline{first<Over1000>}, a fold of one step over the first \cppinline{budget} naturals, the state a candidate and whether it has the property, frozen once it has, read out as an optional, here behind a non-template wrapper.
The fold over \cppinline{iota(0, budget)} with that step is one loop, and its header has three $\phi$-nodes: the two components of the state, the candidate and whether it has been found, and the position in the range.
Those are the iterator's state in Figure~\ref{fig:fjlt-translation}, the pair $(\llbracket r \rrbracket, \llbracket a \rrbracket)$ with $\llbracket a \rrbracket$ itself a pair, and the count $|r|$, here $\llbracket j \rrbracket - \llbracket i \rrbracket$, is the comparison that ends the loop.
The exit block's two $\phi$-nodes are the fold's result, read out into the optional.
Nothing else survives: not the predicate's concept, not the record its constraint became, not the range object, not the pair; the predicate is the multiplication and the comparison in the loop body, and the budget is the register the wrapper was called with.

\begin{listing}[htbp]
\begin{minted}{llvm}
define { i64, i8 } @first_rt(i64 %0) {
  %2 = icmp ult i64 %0, 2
  br i1 %2, label %15, label %3

3:                                                ; loop header
  %4 = phi i8 [ %12, %3 ], [ 0, %1 ]
  %5 = phi i64 [ %10, %3 ], [ 0, %1 ]
  %6 = phi i64 [ %13, %3 ], [ 1, %1 ]
  %7 = trunc i8 %4 to i1
  %8 = mul i64 %6, %6
  %9 = icmp ugt i64 %8, 1000
  %10 = select i1 %7, i64 %5, i64 %6
  %11 = zext i1 %9 to i8
  %12 = or i8 %4, %11
  %13 = add i64 %6, 1
  %14 = icmp eq i64 %13, %0
  br i1 %14, label %15, label %3

15:                                               ; exit
  %16 = phi i64 [ 0, %1 ], [ %10, %3 ]
  %17 = phi i8 [ 0, %1 ], [ %12, %3 ]
  %18 = trunc i8 %17 to i1
  %19 = select i1 %18, i64 %16, i64 0
  %20 = insertvalue { i64, i8 } poison, i64 %19, 0
  %21 = insertvalue { i64, i8 } %20, i8 %17, 1
  ret { i64, i8 } %21
}
\end{minted}
\caption{\cppinline{first<Over1000>(budget)} in LLVM IR, Clang 23 at \texttt{-std=c++23 -O2} with unrolling and vectorisation off, attributes and metadata stripped; Appendix~\ref{app:fjlt:ssa} reads it.}
\label{lst:ssa}
\end{listing}

This is a reason, not a theorem: the compiler is not among the hypotheses, its lowering is trusted, as every \Cpp{} program trusts it, and the reading of SSA as iterator terms is Appel's, which this paper uses and does not prove.
Since both targets make the loop's state a parameter, the sentence of Section~\ref{sec:binding}, that the discipline costs nothing at run time, has a structural ground: a Jlt fold names in the source the parameter the compiler would otherwise recover from assignments.
One seam shows here too: SSA's values are the scalars in registers, and a value that lives in memory, a boxed value in particular, is read and written through a pointer that no $\phi$-node carries, which is the same line between a value and a cell that Rule~1 draws.

